% - Why is asset allocation important? Short description of the asset allocation optimisation problem.
% - Short history of Markowitz's Modern Portfolio Theory.

Because no investor has unlimited capital to invest, all of them are faced with the same challenge: how to allocate the limited resources that they have across the set of investable assets available. Moreover, this task gets more complicated when investors have to take into account their risk appetite and their target return. We can then define the asset allocation problem as an optimisation problem with 2 variants:
\begin{enumerate}
    \item Maximise returns while maintaining risk within a given risk tolerance level, or
    \item Minimise risk for a given level of target return.
\end{enumerate}

Harry M. Markowitz was the first to develop and present a systematic approach to this asset allocation problem, in both variants, in 1952, earning him a Nobel Memorial Prize in Economic Sciences in 1990. Markowitz's portfolio construction model laid the ground for quantitative asset allocation models. Many asset allocation models are developed using Markowitz's model and adjusting it to account for modern investors' needs by for example using other risk measures or more specific target functions other than expected returns or their standard deviation.

The main critique point of Markowitz's portfolio construction model is the lack of stability that leads to concentrated portfolios when correlations rise and to a high portfolio turn-over. In this paper, we analyse Markowitz's portfolio construction model and the causes of low stability. Moreover, we present some techniques that try to improve portfolio stability: 1 regarding parameter estimation for Markowitz's model and 3 alternative allocation methods.

\subsection{Practical Implementation Framework}
    